\chapter{La communication hormonale}\label{ch:g8:hormonal-communication}

La \omterm{def:g8:puberty:puberty}{puberté} les a présentées
(\cref{def:g8:puberty:hormones})~; la pilule s'en sert
(\cref{def:g8:contraception:pill})~; la vague chaude de la
frayeur en a livré une en plein sprint
(\cref{ex:g8:nervous-communication:cases}). Le service
postal mérite sa propre visite~: les \omterm{def:g8:hormonal-communication:glands}{glandes} qui écrivent
les lettres, le \omterm{prop:g4:journey-of-food:absorb}{sang} qui les porte, l'astuce d'adressage qui
rend précise une diffusion générale --- et les boucles de
\omterm{prop:g8:hormonal-communication:feedback}{rétrocontrôle} qui tiennent toute la correspondance en équilibre.

\section{Glandes, hormones, organes cibles}

\begin{definition}[Les glandes endocrines]\label{def:g8:hormonal-communication:glands}
Une \emph{glande}\index{glande} dont le produit est déversé
\emph{dans le \omterm{prop:g4:journey-of-food:absorb}{sang}} (et non le long d'un canal, comme celui
de la \omterm{def:g4:journey-of-food:tract}{salive}) est une glande \emph{endocrine}, et ses
produits sont les \emph{hormones}\index{hormone}. Les
principales rédactrices~: le \emph{centre de commande} à la
base du \omterm{def:g5:brain-in-command:system}{cerveau} --- le starter de la \omterm{def:g8:puberty:puberty}{puberté}, et le chef
d'orchestre de la plupart des autres~; la \emph{thyroïde}
dans le cou --- son hormone règle l'allure de tout le
corps~; les capuchons \emph{surrénaux} posés sur les \omterm{def:g7:removing-wastes:kidneys}{reins}
--- l'hormone d'alarme des frayeurs et des sprints~; le
\emph{pancréas} --- son hormone gère le \omterm{prop:g7:organs-at-work:bill}{glucose} du \omterm{prop:g4:journey-of-food:absorb}{sang}~;
et les \emph{\omterm{def:g8:reproductive-systems:female}{ovaires} et les \omterm{def:g8:reproductive-systems:male}{testicules}} --- les hormones du
cycle et de la reconstruction de
\cref{ch:g8:reproductive-systems}.
\end{definition}

\begin{proposition}[L'adressage par récepteur]\label{prop:g8:hormonal-communication:address}
Le \omterm{prop:g4:journey-of-food:absorb}{sang} livre chaque \omterm{def:g8:hormonal-communication:glands}{hormone} partout
(\cref{ex:g8:puberty:sites}) --- et pourtant chacune n'agit
que sur ses \emph{organes cibles}\index{organe cible}~: les
\omterm{def:g5:first-look-at-cells:cell}{cellules} bâties avec des sites d'amarrage, des
\emph{récepteurs}, pour ce messager. Transport en diffusion
générale, adressage côté destinataire~: une lettre que
toutes les maisons reçoivent, lisible seulement là où la
clé convient. (L'amarrage de la \omterm{def:g8:nervous-communication:synapse}{synapse}, dans
\cref{def:g8:nervous-communication:synapse}, est la même
astuce à l'échelle du nanomètre.)
\end{proposition}
\admitted

\begin{example}[La tournée d'un messager~: l'hormone d'alarme]\label{ex:g8:hormonal-communication:alarm}
Un accident évité de justesse à vélo~: les \omterm{def:g8:hormonal-communication:glands}{glandes}
surrénales déversent leur \omterm{def:g8:hormonal-communication:glands}{hormone} d'alarme dans le \omterm{prop:g4:journey-of-food:absorb}{sang} et,
en quelques secondes --- un tour de circulation
(\cref{ex:g4:heart-and-blood:round}) --- chaque récepteur
équipé répond d'un coup~: le \omterm{def:g4:heart-and-blood:heart}{cœur} accélère, les voies
aériennes s'élargissent, le \omterm{prop:g7:digestion-and-nutrients:absorption}{foie} relâche le \omterm{prop:g7:organs-at-work:bill}{glucose} mis en
réserve (la banque de
\cref{ex:g7:digestion-and-nutrients:breadtrace}, appelée),
les vaisseaux des \omterm{def:g3:bones-and-muscles:muscle}{muscles} s'élargissent et ceux du \omterm{def:g4:journey-of-food:tract}{tube
digestif} se resserrent (le déroutage de
\cref{prop:g7:organs-at-work:supply}, tenu par l'\omterm{def:g8:hormonal-communication:glands}{hormone}),
et la \omterm{def:g1:the-five-senses:sense}{peau} devient pâle et picotante. Une lettre, six
réponses, un seul état coordonné~: prêt.
\end{example}

\section{L'équilibre par rétrocontrôle}

\begin{proposition}[Le principe du rétrocontrôle]\label{prop:g8:hormonal-communication:feedback}
Les systèmes hormonaux se règlent eux-mêmes par
\emph{rétrocontrôle}\index{rétrocontrôle}~: la \omterm{def:g8:hormonal-communication:glands}{glande} qui
écrit \emph{surveille le résultat} de ses propres lettres,
et ralentit quand le résultat est atteint --- un thermostat,
pas un robinet. La vitrine, c'est le \omterm{prop:g7:organs-at-work:bill}{glucose} sanguin~: après
un repas, le \omterm{prop:g7:organs-at-work:bill}{glucose} qui monte déclenche lui-même l'\omterm{def:g8:hormonal-communication:glands}{hormone}
de mise en réserve du pancréas~; le \omterm{prop:g7:organs-at-work:bill}{glucose} rangé, le
niveau baisse, et la sécrétion retombe avec lui~; entre les
repas, une \omterm{def:g8:hormonal-communication:glands}{hormone} partenaire relâche la banque dans
l'autre sens. Le niveau roule dans un couloir étroit, jour
et nuit, sans aucun conducteur conscient.
\end{proposition}
\admitted

\begin{example}[Quand le rétrocontrôle échoue~: le diabète]\label{ex:g8:hormonal-communication:diabetes}
Dans le \emph{diabète}\index{diabète}, le signal de mise en
réserve échoue --- le pancréas cesse de l'écrire, ou les
destinataires cessent de répondre. Le \omterm{prop:g7:organs-at-work:bill}{glucose} s'entasse
dans le \omterm{prop:g4:journey-of-food:absorb}{sang} après chaque repas~; la réabsorption par le
\omterm{def:g7:removing-wastes:kidneys}{rein} est débordée et du sucre apparaît dans l'\omterm{def:g7:removing-wastes:kidneys}{urine} ---
le signal classique de
\cref{ex:g7:removing-wastes:thrift}, enfin expliqué. Le
traitement est ce chapitre mis en pratique~: des doses
mesurées de l'\omterm{def:g8:hormonal-communication:glands}{hormone} manquante, calées sur les repas ---
la lettre écrite à la seringue. Des millions de personnes
vivent très bien de cet arrangement exact.
\end{example}

\begin{remark}[Les deux systèmes en association]\label{rem:g8:hormonal-communication:partnership}
Le télégraphe et la poste de
\cref{prop:g8:nervous-communication:compare} sont une seule
administration. Le centre de commande à la base du \omterm{def:g5:brain-in-command:system}{cerveau}
en est l'interface~: le \omterm{def:g5:brain-in-command:system}{système nerveux} au-dessus, les
\omterm{def:g8:hormonal-communication:glands}{hormones} au-dessous --- la frayeur du cycliste est partie
par les \omterm{def:g5:brain-in-command:system}{nerfs} et a été entretenue par une \omterm{def:g8:hormonal-communication:glands}{hormone}~; la
\omterm{def:g8:puberty:puberty}{puberté} se déroule au rythme lent des \omterm{def:g8:hormonal-communication:glands}{hormones} sous un chef
d'orchestre cérébral~; les brebis de
\cref{ex:g8:reproduction-and-environment:lambs} changent la
longueur du jour (une perception nerveuse) en \omterm{def:g8:hormonal-communication:glands}{hormones} de
reproduction. Des fils rapides là où la vitesse compte, une
poste lente là où le programme doit tourner des mois, et des
bureaux de traduction entre les deux.
\end{remark}

\section{Lire un système hormonal}

\begin{method}[Les quatre questions de toute hormone]\label{met:g8:hormonal-communication:read}
Pour toute \omterm{def:g8:hormonal-communication:glands}{hormone} rencontrée, établis~:
\begin{enumerate}
  \item la \emph{rédactrice}~: quelle \omterm{def:g8:hormonal-communication:glands}{glande}, dans le \omterm{prop:g4:journey-of-food:absorb}{sang},
        sur quel déclencheur~?
  \item les \emph{destinataires}~: quels \omterm{prop:g8:hormonal-communication:address}{organes cibles}
        portent ses récepteurs, et que fait chacun à la
        livraison~?
  \item le \emph{\omterm{prop:g8:hormonal-communication:feedback}{rétrocontrôle}}~: quel résultat fait son
        rapport, et comment l'excès d'écriture est-il évité~?
  \item la \emph{panne et l'usage}~: que se passe-t-il quand
        la lettre manque --- et la médecine l'imite-t-elle
        (la pilule, les doses du diabétique)~?
\end{enumerate}
\end{method}

\begin{example}[La méthode sur la meneuse d'allure thyroïdienne]\label{ex:g8:hormonal-communication:thyroid}
Rédactrice~: la thyroïde, dirigée par le centre de commande.
Destinataires~: presque toutes les \omterm{def:g5:first-look-at-cells:cell}{cellules} --- l'\omterm{def:g8:hormonal-communication:glands}{hormone}
règle la vitesse du feu lent de
\cref{prop:g7:what-is-respiration:power}. \omterm{prop:g8:hormonal-communication:feedback}{Rétrocontrôle}~: le
centre de commande lit le niveau et ajuste sa direction.
Panne~: trop peu, et tout le moteur tourne au ralenti ---
fatigue, froid, tout ralenti~; trop, et il s'emballe. Usage~:
l'\omterm{def:g8:hormonal-communication:glands}{hormone} manquante est un comprimé quotidien --- l'une des
plus anciennes lettres par la poste de la médecine.
\end{example}

\section{Exercices}

\begin{exercise}[$\star$]\label{exo:g8:hormonal-communication:1}
Qu'est-ce qui rend une \omterm{def:g8:hormonal-communication:glands}{glande} endocrine~? Nommes-en quatre,
avec un rôle d'\omterm{def:g8:hormonal-communication:glands}{hormone} pour chacune.
\end{exercise}

\begin{exercise}[$\star$]\label{exo:g8:hormonal-communication:2}
Le \omterm{prop:g4:journey-of-food:absorb}{sang} livre partout~; les \omterm{def:g8:hormonal-communication:glands}{hormones} agissent quelque part.
Résous le paradoxe.
\end{exercise}

\begin{exercise}[$\star$]\label{exo:g8:hormonal-communication:3}
Cite quatre des six réponses à l'\omterm{def:g8:hormonal-communication:glands}{hormone} d'alarme, chacune
avec le chapitre qui la sert.
\end{exercise}

\begin{exercise}[$\star$]\label{exo:g8:hormonal-communication:4}
Énonce le principe du \omterm{prop:g8:hormonal-communication:feedback}{rétrocontrôle}, et sa vitrine.
\end{exercise}

\begin{exercise}[$\star$]\label{exo:g8:hormonal-communication:5}
Qu'est-ce qui échoue dans le \omterm{ex:g8:hormonal-communication:diabetes}{diabète}, et qu'est-ce qui
réapparaît au \omterm{def:g7:removing-wastes:kidneys}{rein} à cause de cela~?
\end{exercise}

\begin{exercise}[$\star$]\label{exo:g8:hormonal-communication:6}
Où les \omterm{def:g5:brain-in-command:system}{systèmes nerveux} et hormonal se rejoignent-ils~?
Donne un exemple de relais.
\end{exercise}

\begin{exercise}[$\star\star$]\label{exo:g8:hormonal-communication:7}
Pourquoi le thermostat est-il la bonne image du
\omterm{prop:g8:hormonal-communication:feedback}{rétrocontrôle}, et le robinet la mauvaise~?
\end{exercise}

\begin{exercise}[$\star\star$]\label{exo:g8:hormonal-communication:8}
Applique \cref{met:g8:hormonal-communication:read} à
l'\omterm{def:g8:hormonal-communication:glands}{hormone} de mise en réserve du pancréas.
\end{exercise}

\begin{exercise}[$\star\star$]\label{exo:g8:hormonal-communication:9}
Explique la seringue du diabétique et la \omterm{def:g8:contraception:pill}{pilule
contraceptive} comme un même genre de geste, avec la
question~4 de la méthode.
\end{exercise}

\begin{exercise}[$\star\star$]\label{exo:g8:hormonal-communication:10}
Attribue les systèmes, avec tes raisons~: le début du
frisson~; la pousse de la barbe~; le couloir étroit du
\omterm{prop:g7:organs-at-work:bill}{glucose} sanguin~; le réflexe rotulien.
\end{exercise}

\begin{exercise}[$\star\star$]\label{exo:g8:hormonal-communication:11}
Les brebis de \cref{ex:g8:reproduction-and-environment:lambs}
changent la \omterm{def:g6:exploring-our-environment:conditions}{lumière} en agneaux. Suis le relais à travers les
bureaux de
\cref{rem:g8:hormonal-communication:partnership}.
\end{exercise}

\begin{exercise}[$\star\star\star$]\label{exo:g8:hormonal-communication:12}
«~Les récepteurs font d'une diffusion générale un réseau de
lignes privées.~» Défends la phrase, puis pousse-la plus
loin~: que doit changer dans une \omterm{def:g5:first-look-at-cells:cell}{cellule} pour qu'elle
\emph{entre} dans le public d'une \omterm{def:g8:hormonal-communication:glands}{hormone} --- et qu'est-ce
que cela suggère sur la façon dont les cibles de la \omterm{def:g8:puberty:puberty}{puberté}
se sont mises à l'écoute~?
\end{exercise}

\section{Problème~: le carnet du glucose}

\begin{problem}[{Problème du week-end --- une journée de
sucre sanguin, lue comme une correspondance}]\label{pb:g8:hormonal-communication:1}
Le relevé (inventé) d'une journée sur un capteur de \omterm{prop:g7:organs-at-work:bill}{glucose}
en continu~: stable la nuit~; petit-déjeuner à 8~h --- une
montée, puis une descente calme jusqu'à 10~h~; sport à
12~h --- un bref creux, vite rattrapé~; un déjeuner sauté
--- stable tout l'après-midi malgré tout~; une frayeur à
17~h (un accident évité de justesse à vélo) --- une montée
brève et brusque~; dîner à 19~h --- montée et descente à
nouveau.

\medskip\noindent\textbf{Partie I --- La matinée.}
\begin{enumerate}
  \item La stabilité nocturne~: quelle correspondance à
        deux lettres a tenu le couloir pendant que son
        propriétaire dormait~?
  \item La montée et la descente du petit-déjeuner~: nomme
        le déclencheur, la rédactrice, la lettre et les
        destinataires qui ont mis le surplus en réserve.
  \item Pourquoi \emph{l'arrêt} de la descente --- aucun
        plongeon sous le couloir --- est-il à lui seul la
        signature du \omterm{prop:g8:hormonal-communication:feedback}{rétrocontrôle}~?
  \item Le niveau de 10~h égale celui de 7~h. Que prédit
        l'image du thermostat sur la courbe de la sécrétion
        entre les deux~?
\end{enumerate}

\medskip\noindent\textbf{Partie II --- L'après-midi.}
\begin{enumerate}[resume]
  \item Le creux du sport et son rattrapage~: quelle lettre
        a appelé la banque du \omterm{prop:g7:digestion-and-nutrients:absorption}{foie}, et le compte de quel
        chapitre était-il en train d'être dépensé~?
  \item Le déjeuner sauté n'a rien changé de visible. Le
        travail discret de quelle \omterm{def:g8:hormonal-communication:glands}{hormone} partenaire a-t-il
        fait le pont de l'après-midi~?
  \item La montée brusque de la frayeur~: nomme la
        rédactrice et le sens de ce sucre soudain --- prêt à
        quoi, d'après
        \cref{ex:g8:hormonal-communication:alarm}~?
  \item Pourquoi la montée de la frayeur est-elle retombée
        dans l'heure, contrairement à l'entassement durable
        du \omterm{ex:g8:hormonal-communication:diabetes}{diabète}~?
\end{enumerate}

\medskip\noindent\textbf{Partie III --- La consultation.}
\begin{enumerate}[resume]
  \item Le carnet d'un diabétique, pour la même journée,
        différerait à chaque repas. Décris la courbe du
        petit-déjeuner sans traitement, et dis où
        apparaîtrait le signal du \omterm{def:g7:removing-wastes:kidneys}{rein}.
  \item La version traitée~: qu'imite la dose mesurée prise
        avant chaque repas, et pourquoi sa taille doit-elle
        correspondre à l'assiette~?
  \item Le médecin appelle le carnet sain «~une conversation
        que l'on n'entend jamais~». Dresse la liste des
        interlocuteurs de la journée --- \omterm{def:g8:hormonal-communication:glands}{glandes}, lettres,
        destinataires --- en un tableau ou un paragraphe de
        synthèse.
  \item Termine par la morale du chapitre en une phrase~:
        quel genre de système tient une valeur dans un
        couloir toute une vie sans conducteur~?
\end{enumerate}
\end{problem}
